\section{Introduction}
\label{sec:intro}

Anomaly detection in videos is a critical task with applications across many domains such as surveillance, healthcare, and autonomous driving\citep{a-liu2021, a-garg2020}. The goal is to identify irregular events that deviate from normal patterns in video data. However, annotated real-world video datasets containing anomalies are limited due to the prohibitive cost of exhaustive manual annotation. This has led to increasing interest in weakly supervised anomaly detection\citep{a-jiang2023weakly}, where models are trained on videos with only video-level labels indicating the presence of anomalies instead of frame-level or pixel-level annotations. 

\begin{figure}[t]
\centering
\includegraphics[width=\linewidth]{images/arch.png}
\caption{The overall framework of the proposed method. The blue line is the data flow of the classical method and the red line is the data flow of our method. The three modules marked with stars are the novelties of our approach.}
\label{fig:arch}
\end{figure}

While weakly supervised methods reduce annotation costs, modeling normal activities and detecting anomalies from only video-level labels remains challenging. Existing approaches\citep{a-wu2022self,a-li2022self,a-tian2021weakly,a-sultani2018real} have focused on extracting discriminative features from individual frames and aggregating them to perform video-level classification. However, these methods do not explicitly model the temporal evolution of activities and often struggle to detect anomalies that persist over long durations. Recent work\citep{a-shin20203d,a-murugesan2020efficient} has incorporated temporal modeling using 3D ConvNets or recurrent neural networks (RNNs) over frame-level features. However, performance is still limited because these models lack sufficient supervision and tend to overfit to frequent normal activities while failing to identify infrequent anomalies effectively. Methods based solely on feature recognition or temporal information have reached a performance bottleneck. Therefore, as shown in Figure \ref{fig:arch}], we hypothesize that extracting richer information from videos and establishing prior knowledge for the targeted anomaly types, and then incorporating this knowledge as auxiliary supervision during training, would further improve the performance and generalization capability of the model.

In this work, we aim to improve weakly supervised anomaly detection by incorporating external knowledge from textual descriptions to provide richer supervision. Specifically, we leverage recent advances in large language models (LLMs) to generate textual prompts that describe normal activities and anomalies for a given dataset\citep{a-zhao2023survey}. These prompts prime the model with high-level semantic information about the appearance and evolution of events beyond what is contained in the video-level labels. 

However, directly incorporating raw prompt embeddings as additional input can be inefficient and redundant. Hence, we propose to first normalize the prompt embeddings using dictionary learning techniques\citep{a-tovsic2011dictionary} like Sparse Coding\citep{a-lee2006efficient} and Dense Associative Memory\citep{a-krotov2016dense}. This allows compressing the high-dimensional prompt embeddings into a compact dictionary of codewords that capture the essential semantic concepts. The codeword activations for a given prompt serve as its normalized representation. 

More formally, let $p_i \in \mathbb{R}^d$ be the embedding for prompt $i$ obtained from the LLM, where $d$ is the dimensionality. We learn a dictionary $D \in \mathbb{R}^{d \times K}$ with $K$ codewords that sparsely represents the prompt embeddings:

\begin{equation}
p_i \approx D a_i, \quad ||a_i||_0 \ll K
\end{equation}

where $a_i \in \mathbb{R}^K$ is the sparse code for prompt $i$ containing mostly zero activations, except for a few codewords relevant to the semantic content. We can compute the code activations efficiently using sparse coding algorithms like orthogonal matching pursuit. The code $a_i$ serves as the normalized prompt representation that will be provided as auxiliary input to the detection model.

Furthermore, during our experiments, we discovered a significant inconsistency issue in the scene text descriptions extracted from video frames, leading to ambiguity and distortion. This directly impacts the algorithm's ability to comprehend and recognize video content. To address this, we propose a novel consistency regularization method. For each video, we first obtain frame-level image captions using a pre-trained image captioning model such as BLIP\citep{a-li2022blip,a-li2023blip}. We then enforce consistency between frames with high similarity in the video by further using the LLM to generate the description from the original captions and extra supporting information generated by other models. Concretely, if $c_{i,1}, c_{i,2}, ... c_{i,T}$ are the frame-level captions for video $i$ and $d_{i,1}, d_{i,2}, ... d_{i,T}$ are extra supporting information gained by other models such as YOLO, we optimize: 

\begin{equation}
 c_{d_{i, 1}}, ..., c_{d_{i, T}} = \text{LLM}_\theta(c_{i,1}, ..., c_{i,T}, d_{i,1}, ..., d_{i,T})
\end{equation}

where $c_{d_{i, T}}$ is the textual description after LLM correction and $\theta$ denotes the LLM parameters. This regularization method corrects erroneous raw text descriptions with additional information.

For modeling, we propose a novel cross-modal architecture that leverages the normalized prompt embeddings. Our model consists of three branches: an RGB stream, an optical flow stream\citep{a-fortun2015optical}, and a text stream. The RGB and flow streams extract spatio-temporal features from frame sequences using pre-trained visual models\citep{a-carreira2017quo}. The text stream encodes the prompt representation $a_i$ and $c_{d_{i, T}}$ through an embedding layer followed by modeling with BART\citep{a-lewis2019bart}. 

We introduce cross-modal attention mechanisms that allow bi-directional interaction between the modalities. The text stream attends to relevant normal and anomalous patterns in the video streams using a text-visual attention. Conversely, the video streams focus on regions described in the text using a visual-text attention. This multi-stage attention exchange enables deeply fused integration of the complementary modalities. After that, the prior knowledge further optimizes the fused features by increasing or suppressing the expressiveness of different regions.

Finally, the fused features are passed through classification layers to predict frame-level anomalies. The cross-modal architecture is trained end-to-end using the video labels only. We empirically demonstrate significant gains over state-of-the-art approaches on multiple benchmark datasets. 

In summary, our key contributions are:

\begin{itemize}
    \item Leveraging LLMs to generate informative textual prompts as external knowledge and normalizing textual prompts via dictionary learning for efficient modeling
    \item Enforcing caption consistency using LLMs and extra supporting information
    \item Cross-modal neural architecture with bi-directional attentive fusion
    \item Demonstrating improved anomaly detection on two benchmark datasets
\end{itemize}

